\begin{abstract}
Systems for multi-event video retrieval are usually reported with one competition score, which hides two failure modes: the first pass returns the wrong video, or it returns the right video but places some events at the wrong moments.
We study both on the Ho Chi Minh City AI Challenge corpus (873 videos, 470{,}804 frames) with a pipeline that fuses SigLIP2, BGE-M3, and BM25 evidence and aligns ordered events with dynamic programming (DP).
We annotate per-event intervals for 40 queries (113 events) and video-level answers for 70 queries.
Even when the correct video is given, only 25\% of queries have every event placed correctly, and a correct frame is ranked first for only 35\% of events.
The relative score margin between the top video and the next distinct video separates correct from wrong first passes with AUROC 0.85 overall, and nested leave-one-set-out signal selection gives held-out AUROCs of 0.78--0.85.
With the correct video given, simulated oracle feedback that keeps, uses, or rejects one event at a time raises whole-sequence correctness from 25\% to 70\% within five actions with eight alternatives per event, against 47.5\% without DP; the gain comes mostly from the DP starting path and order constraints rather than from propagation to untouched events.
These results suggest substantial headroom for failure-aware interaction without retraining, and identify within-video frame scoring as the primary bottleneck in the studied pipeline.
\keywords{Multi-Event Video Retrieval \and Temporal Alignment \and Query Performance Prediction \and Interactive Correction}
\end{abstract}
